\chapter{Kết luận và Hướng phát triển}
\label{ch:conclusion}

Chương này tổng kết kết quả nghiên cứu và đóng góp của luận văn, khẳng định giá trị khoa học lẫn thực tiễn, phân tích các giới hạn kiến trúc và phác thảo lộ trình phát triển.

\section{Tổng hợp Kết quả Nghiên cứu}
\label{sec:conclusion_summary}

\textbf{Về RQ1}, các phép đo xác nhận kiến trúc phân tầng giải quyết được vấn đề tiêu thụ tài nguyên của mô hình ngôn ngữ. Trên luồng benchmark 25.649 sự kiện với 31,6\% tấn công, Tầng 1 và Cổng Học máy giải quyết \textbf{90,6\%} lưu lượng mà không phát sinh lời gọi LLM nào; trên luồng 99.717 sự kiện ở tỉ lệ tấn công gần thực tế hơn (9,8\%), con số này đạt \textbf{97,5\%}. Nghiên cứu chỉ ra tỉ lệ xả tải là \textbf{hàm của tỉ lệ tấn công trong luồng} chứ không phải hằng số của hệ: trong mỗi luồng, tỉ lệ tổng tách tuyến tính thành một tỉ lệ cho ca lành và một tỉ lệ cho ca tấn công---92,35\% và 86,72\% trên luồng benchmark, 97,88\% và 93,77\% trên luồng dạng SOC---nên quan hệ này cho phép chiếu sang tỉ lệ tấn công khác trong cùng một luồng, chứ không cho phép mang hệ số của luồng này áp sang luồng kia. Cổng ML xử lý mỗi véc-tơ trong 0,287\,ms (3.452 sự kiện/giây), bộ đệm ngữ nghĩa đạt tỉ lệ trúng 81,3\% với hệ số tăng tốc 8,9$\times$, và vòng phản hồi từ Tầng 2 về Tầng 1 giảm thêm 11,35\% tải leo thang \textit{đồng thời} cải thiện chất lượng phát hiện (MCC $+0{,}0378$). Toàn tuyến, độ trễ trung bình giảm từ \textbf{17.187,9\,ms} xuống \textbf{5.286,7\,ms} ($-69{,}2\%$, Mann-Whitney $p \approx 9 \times 10^{-57}$) và trung vị từ 17.174,7\,ms xuống \textbf{0,88\,ms}. Luận văn báo cáo ngang bằng kết quả ngược chiều: phân vị 95 \textbf{xấu hơn} (25.829\,ms so với 21.434\,ms), vì sự kiện nào leo thang thì đã trả chi phí Tầng 1 và Cổng ML trước khi trả chi phí LLM. Lợi ích tập trung ở trung vị chứ không trải đều.

\textbf{Về RQ2}, chuỗi niêm phong HMAC-SHA256 phát hiện \textbf{100\%} ba kiểu giả mạo cốt lõi---sửa nội dung, chèn bản ghi giả và xoá bản ghi giữa---trên 30 lượt thử mỗi kiểu, với đối chứng âm trên 450 bản ghi nguyên vẹn không sinh báo động giả. Phạm vi bảo chứng được phát biểu chính xác: vì dùng khóa đối xứng, cơ chế bảo đảm \textbf{toàn vẹn và phát hiện giả mạo}, không phải chống chối bỏ theo nghĩa mật mã học; và \textbf{cắt cụt đuôi chuỗi là điểm mù nguyên lý}. Hệ thống được đánh giá trước hai mặt tấn công khác biệt. Trước tấn công \textbf{qua ngôn ngữ}, lớp lọc tĩnh chặn 18,8\% trên 703 mẫu từ nguồn công bố---kèm tỉ lệ báo nhầm 20,0\% trên log lành mà toàn bộ tập trung ở lớp nhận dạng mã hoá---trong khi Tầng 2, đo ở \textbf{độ phủ 100\%} của 678 mẫu thuộc các nhóm mà lớp tĩnh không hấp thụ được, kháng được \textbf{678/678}. Trước tấn công \textbf{qua không gian đặc trưng}, nơi giá trị cực đoan được bơm vào để ép Cổng ML nói ``lành'', tỉ lệ kháng đạt \textbf{98,75\%} ở chế độ khó bóp méo toàn bộ đặc trưng (1.023/1.036; CI95 97,86--99,27). Kết quả này củng cố luận điểm trung tâm: bảo chứng an ninh đến từ \textbf{Đóng gói Dữ liệu Phân định}, vốn xử lý nội dung trong vùng bọc như văn bản thụ động bất kể nó nói gì, nên không phụ thuộc vào năng lực nhận dạng nội dung độc hại của bất kỳ bộ phân loại nào. Phép kiểm độ ổn định xác nhận tác tử tất định với cùng seed và không đổi phán quyết nào khi chỉ đổi seed, nên chênh lệch đo được giữa các cấu hình không phải nhiễu của bộ giải mã.

\textbf{Về RQ3}, tác tử LangGraph Tầng 2 kết hợp Dual-RAG trên 433 mục ATT\&CK chuẩn STIX, được trọng tài khác họ mô hình (Meta-Llama-3-8B) chấm trên 277 sự kiện leo thang, đạt trung bình \textbf{3,73/5,0} với độ liên quan 4,22 và độ trung thành 3,88, trong khi \textbf{độ chính xác ngữ cảnh 2,68 là điểm yếu rõ ràng} chỉ ra nút thắt nằm ở tầng truy xuất; tỉ lệ neo bằng chứng đạt 50,2\%. Đóng góp kiến trúc ở RQ3 là \textbf{rào chắn neo-RAG tất định}: mọi mã kỹ thuật công bố buộc phải hiện diện trong tài liệu truy xuất của chính lô đó, nếu không sẽ bị ép về \texttt{N/A} và chuyển chuyên gia---biến việc chống hư cấu từ kỳ vọng đặt vào mô hình thành ràng buộc kiểm chứng được, giữ nguyên hiệu lực khi đổi mô hình. Luận văn \textbf{không công bố} tỉ lệ khớp chính xác mã ATT\&CK ở đợt này, vì các tệp kết quả hiện có được sinh trước giao ước loại mẫu biên soạn và mâu thuẫn lẫn nhau.

\section{Đóng góp Khoa học và Thực tiễn}
\label{sec:contributions}

Về \textbf{đóng góp khoa học}, luận văn đề xuất và định lượng nguyên lý \textit{đặt phán quyết rẻ trước phán quyết đắt}, giải quyết rào cản độ trễ bậc hai $\mathcal{O}(N^2)$ của mô hình ngôn ngữ trên luồng dữ liệu liên tục~\cite{vaswani2017}. Đóng góp phương pháp luận đi kèm là việc phát biểu tường minh rằng \textbf{tỉ lệ xả tải là hàm của base-rate}, kèm quan hệ tuyến tính cho phép chuyển kết quả sang môi trường có tỉ lệ tấn công khác---điều mà cách báo cáo một con số duy nhất không làm được. Thứ hai, luận văn đóng góp mô hình cô lập dữ liệu bằng Nonce mật mã, cung cấp bảo chứng an ninh \textbf{mang tính cấu trúc}, độc lập với các bộ phân loại xác suất vốn suy giảm trước biến thể mới~\cite{promptinjection2023, owasp_llm2025}. Thứ ba, luận văn xác lập bộ giao ước đo lường cho hệ tác tử an ninh cục bộ: mẫu số khớp câu hỏi, thước đo miễn nhiễm lệch lớp, loại trừ mẫu tự biên soạn khỏi tỉ lệ công bố, và bắt buộc báo tỉ lệ chặn kèm tỉ lệ báo nhầm.

Về \textbf{đóng góp thực tiễn}, luận văn cung cấp nguyên mẫu hoàn chỉnh đóng gói bằng Docker Compose, kiểm chứng trên phần cứng thương mại phổ thông (Intel Core i7, 32\,GB RAM, một card RTX 4060 Ti 16\,GB VRAM), tuân thủ nguyên tắc Zero-Trust trong môi trường air-gapped~\cite{nist80207}. Hệ thống tích hợp bảng điều khiển Streamlit kết nối hàng đợi Con-người-trong-vòng-lặp nhằm thay thế kịch bản SOAR tĩnh~\cite{soar2025}, kèm niêm phong vết pháp y bằng chuỗi băm HMAC-SHA256~\cite{hmac1997}.

\section{Giới hạn và Hướng Phát triển}
\label{sec:limitations}

Bốn giới hạn dưới đây được nêu kèm giải pháp đối ứng tương ứng.

\textbf{Thứ nhất---phạm vi nguồn nhật ký.} Đánh giá hiện tập trung vào đặc trưng NetFlow chuẩn hóa và nhật ký HTTP Tầng 7~\cite{ids2018, dapt2020}. Các nguồn phi cấu trúc (Windows Event Log, Syslog thô, AWS CloudTrail) đòi hỏi cơ chế phân tích cú pháp phức tạp hơn. Hướng đối ứng là xây dựng cỗ máy trích xuất tích hợp thuật toán phân nhóm biểu mẫu log như Drain~\cite{he2017drain} để chuẩn hóa mọi định dạng về véc-tơ đặc trưng trước khi vào đường ống.

\textbf{Thứ hai---ranh giới chịu tải trên một nút.} Khi lưu lượng vượt 100.000 sự kiện/giây, lượng sự kiện leo thang lên Tầng 2 trên GPU đơn 16\,GB có thể vượt năng lực suy luận tức thời, gây tích tụ hàng đợi. Hướng đối ứng là chuyển bộ lọc Tầng 1 sang ngôn ngữ biên dịch (Go hoặc Rust), di trú Redis Stream sang Apache Kafka phân mảnh ghi đĩa, và nâng cấp Tầng 2 lên vLLM hoặc TGI hỗ trợ PagedAttention và continuous batching.

\textbf{Thứ ba---quy kết kỹ thuật và chất lượng truy xuất.} Trục độ chính xác ngữ cảnh 2,68 cho thấy tầng truy xuất đưa vào nhiều tài liệu không liên quan; đồng thời lượt đo chất lượng lập luận có 69/277 phán quyết thiếu trường bắt buộc nên chỉ mang tính định hướng. Hướng đối ứng gồm: đo lại quy kết theo giao ước đã chuẩn hoá kèm sàn đoán bừa 52\%, bổ sung chỉ số truy xuất (recall@$k$, MRR) để tách lỗi truy xuất khỏi lỗi sinh, và kiểm định độ tin cậy của trọng tài LLM bằng hệ số Cohen's $\kappa$ đối chiếu với người chấm.

\textbf{Thứ tư---tương quan chuỗi tấn công và tự từ chối dịch vụ.} Cơ chế liên kết hành vi hiện dựa trên địa chỉ IP nguồn cố định, nên trong môi trường xoay IP hoặc botnet phân tán, hiệu quả sẽ giảm. Đồng thời, để tránh chặn nhầm hạ tầng nội bộ, hệ thống buộc các phán quyết dựa trên quy tắc chung chung về hàng đợi chuyên gia, chấp nhận giảm mức tự động hóa. Hướng đối ứng là phát triển tương quan dựa trên \textbf{véc-tơ kỹ thuật} thay vì địa chỉ nguồn, kết hợp Mạng Thần kinh Đồ thị, và áp dụng Học tăng cường từ Phản hồi của Con người trên các quyết định phê duyệt tại hàng đợi \texttt{AWAIT\_HITL} để dần tự động hóa an toàn.

Ngoài bốn hướng trên, lộ trình dài hạn bao gồm nâng chuỗi niêm phong lên chữ ký bất đối xứng (Ed25519) nhằm đạt tính chống chối bỏ thực sự, neo chu kỳ giá trị băm ra ngoài hệ thống để bịt điểm mù cắt đuôi, và tích hợp chuẩn Model Context Protocol cho phép tác tử điều phối động các công cụ an ninh mở rộng.

\vspace{0.3cm}
\begin{flushright}
    \textit{Khép lại một nỗ lực nghiên cứu nghiêm túc, nhưng mở ra những khát vọng sâu xa trong hành trình bảo vệ không gian số bằng sức mạnh của Trí tuệ Nhân tạo Tác tử.}
\end{flushright}
